% ════════════════════════════════════════════════════════════════════
% SECTION 2 — BACKGROUND AND RELATED WORK
% ════════════════════════════════════════════════════════════════════
\section{Background and Related Work}\label{sec:related}

\subsection{Solar flare prediction from photospheric magnetograms}

The physical premise of data-driven flare forecasting is that the
photospheric magnetic field of an active region encodes the stress
budget that a future flare will release. Schrijver~\cite{schrijver2007}
demonstrated that a single scalar---the R-value, the unsigned flux
within high-gradient polarity-separation corridors---separates
major-flaring regions from non-flaring ones with striking reliability,
and this proxy remains among the strongest single flare predictors
identified to date. Bobra and Couvidat~\cite{bobra2015} operationalised
the machine-learning view of the problem, showing that a support-vector
machine trained on 25~SHARP parameters derived from SDO/HMI vector
magnetograms can forecast M-class flares with useful skill, and that
feature selection matters as much as model choice. Subsequent
community-wide blind tests, most notably the multi-method comparison of
Leka et al.~\cite{leka2019}, established two durable findings: (i)~no
single method dominates across all event definitions and skill
metrics, and (ii)~the effective ceiling on skill is set by the quality
and homogeneity of the training data rather than by classifier
complexity.

The European FLARECAST project~\cite{georgoulis2021} pushed the
operational agenda further, combining multiple forecasting methods and
data sources in an open framework, and it explicitly confronted the
data-heterogeneity problem that motivates the present paper. On the
data side, the SWAN-SF benchmark of Angryk et
al.~\cite{angryk2020} released a multivariate time-series dataset of
roughly 8{,}000 active-region tracks with 24 physical parameters per
timestep and flare-history labels, specifically designed to make flare
prediction reproducible and comparable; it has since become the
community's reference benchmark, and a family of cleaned and rebalanced
variants (random undersampling, Tomek-link cleaning, and TimeGAN-based
synthetic augmentation \cite{yoon2019}) has been curated to make the
benchmark tractable for supervised learning. Recent deep-learning work
on SWAN-SF continues to raise the centralised ceiling: contrastive
pretraining and multifaceted preprocessing \cite{eskandari2024} and
LSTM-style temporal architectures \cite{hassani2025} exploit the
sequential structure of the 60-minute observation windows. All of these
works, without exception, assume that the full training corpus is
available in one place.

\subsection{Federated learning under statistical heterogeneity}

Federated learning was introduced by McMahan et
al.~\cite{mcmahan2017} in the FedAvg formulation: clients perform
several epochs of local SGD on private data, and a server aggregates
the resulting weights by sample-size-weighted averaging. The defining
difficulty, identified almost immediately, is statistical
heterogeneity: when client data distributions differ---the non-IID
setting---local optima diverge from the global optimum, and naive
averaging can oscillate or diverge \cite{zhao2018}. FedProx
\cite{li2020} addresses this by adding a proximal term
$(\mu/2)\norm{\w-\w^{(t)}}^{2}$ to each client's local objective that
anchors local training to the current global model, tolerating partial
work and heterogeneous local computation. SCAFFOLD \cite{karimireddy2020}
attacks the same problem with control variates that correct client
drift using variance reduction. The comprehensive survey of Kairouz et
al.~\cite{kairouz2021} situates these algorithms within the broader
open-problem landscape, including robust aggregation, personalisation,
and privacy. Cross-silo FL---few, reliable, institutionally owned
clients---is the deployment regime relevant to space-weather agencies,
and it contrasts with the cross-device regime in its stability
requirements and long-lived participants.

\subsection{Class imbalance in federated learning}

Flare prediction inherits an extreme class-imbalance regime: in
SWAN-SF's natural test distribution, roughly 1.9\% of windows precede
M/X-class flares, while the rebalanced training partitions sit near
parity. The focal loss of Lin et al.~\cite{lin2017}, originally
developed for dense object detection, down-weights easy examples with a
$(1-p_t)^{\gamma}$ modulating factor and concentrates gradient on hard
positives; it has become the default imbalance workhorse. Sarkar et
al.~\cite{sarkar2020} adapted focal loss specifically to FL in the
Fed-Focal formulation, in which the class-weighting parameter is
adjusted using information available at the server about the global
class distribution, so that clients with different local positivity
rates do not implicitly reweight the global objective. In parallel,
synthetic oversampling remains a standard pre-processing response:
SMOTE \cite{chawla2002} interpolates minority-class neighbours, and
mixup \cite{zhang2018} creates convex combinations of training pairs.
Under federation, rebalancing must be applied locally per client, since
the pooled distribution is precisely what no participant can inspect.

\subsection{Positioning of this work and search protocol}
\label{sec:litsearch}

Two research threads---physics-grounded flare prediction and
heterogeneity-robust federated learning---have developed largely
independently, and we are not aware of any prior work that joins them.
Because the first-FL-treatment claim is a factual statement about the
literature, we document the search protocol behind it rather than
relying on ``to the best of our knowledge''. The search was last
repeated on 29~September 2026 across general web, arXiv, NASA ADS,
and publisher indexes (ScienceDirect, IEEE) reached through a web
search engine, using four query strings: (i)~\emph{federated learning
solar flare prediction}; (ii)~\emph{federated learning space weather
forecasting}; (iii)~\emph{federated learning SWAN-SF benchmark
magnetogram}; (iv)~\emph{cross-silo federated learning geophysics
heliophysics distributed observatories}. Records were included if
they apply federated or distributed multi-client learning without raw
data pooling to flare prediction or space-weather forecasting; records
were excluded if they concern centralised machine learning for flares
(the case for all entries of query~i), federated forecasting of
\emph{solar power or irradiance}, federated \emph{terrestrial}
weather forecasting, or federated computing on spacecraft without a
scientific forecasting task (the categories returned by query~ii).
Screening the 30 records returned by the four queries yielded
\emph{zero} in-scope records; the raw result listings are committed
with the repository (\texttt{docs/lit\_search/}). We note two
limitations of this protocol: web-engine coverage of grey literature
is incomplete, and the screening was single-reviewer, so a
PRISMA-style dual-screened log remains queued repository work. Within
those limits, the present work is the first federated treatment of
solar flare prediction, the first FL framework evaluated on the
SWAN-SF benchmark, and the first to quantify the cost of the
data-locality constraint---the gap between federated and centralised
skill---specifically for a safety-critical smart-grid application. In
doing so it also contributes a rarely documented artifact: an honest
engineering log of the failure modes encountered when FL meets a
severely imbalanced physical-science benchmark (Appendix~\ref{sec:appendix}).
